\metatitle{Column formulas for OEIS A358628}
\metadescription{Proofs of the A358628 column formulas, factorization, and Ehrhart--Legendre structure.}


\section[oeisProofA358628]{Column formulas for A358628}

This page is part of the \hyperref[oeisProofs]{OEIS proof collection}.

The square array \oeis{A358628} is defined by
\[
A(i,j)\coloneqq
\sum_{\substack{X,Y\in\setN^j\\ |X|,|Y|\leq i}}
\prod_{k=1}^j(1+X_k+Y_k).
\]
We prove the conjectured generating function for each column and the
conjectured polynomial factorization in the first index.

\begin{theorem}
For $i,j\geq0$,
\[
A(i,j)=[x^iy^i]\frac{(1-xy)^j}
{(1-x)^{2j+1}(1-y)^{2j+1}}.
\]
Consequently, the generating function of column $j$ is
\[
\sum_{i\geq0}A(i,j)t^i
=\frac{\displaystyle\sum_{k=0}^{2j}\binom{2j}{k}^2t^k}
{(1-t)^{3j+1}}.
\]
Equivalently,
\[
\begin{aligned}
A(i,j)
&=\sum_{m=0}^j(-1)^m\binom jm
  \binom{i-m+2j}{2j}^2\\
&=\sum_{k=0}^{2j}\binom{2j}{k}^2
  \binom{i+3j-k}{3j}.
\end{aligned}
\]
\end{theorem}

\begin{proof}
The summand factors over the coordinates, and
\[
\sum_{a,b\geq0}(1+a+b)x^ay^b
=\frac{1-xy}{(1-x)^2(1-y)^2}.
\]
The two bounds $|X|,|Y|\leq i$ contribute the partial-sum factor
$1/((1-x)(1-y))$.  Extracting $[x^iy^i]$ gives the first formula.

Expanding $(1-xy)^j$ and taking the diagonal gives
\[
\sum_{i\geq0}A(i,j)t^i
=(1-t)^j\sum_{n\geq0}\binom{n+2j}{2j}^2t^n
=(1-t)^j\,{}_2F_1(2j+1,2j+1;1;t).
\]
Euler's hypergeometric transformation changes the last expression to
\[
(1-t)^{-3j-1}\,{}_2F_1(-2j,-2j;1;t)
=\frac{\sum_{k=0}^{2j}\binom{2j}{k}^2t^k}
{(1-t)^{3j+1}}.
\]
Coefficient extraction from the two displayed forms gives the alternating
and subtraction-free formulas for $A(i,j)$.
\end{proof}

The subtraction-free formula has an Ehrhart interpretation.  Let
\[
\mathcal P_j\coloneqq
\left\{(x,y,s)\in\setR_{\geq0}^{3j}:
\sum_kx_k\leq1,\quad \sum_ky_k\leq1,\quad
s_k\leq x_k+y_k\text{ for all }k\right\}.
\]
For fixed lattice points $X,Y$ in the two simplices, there are
$\prod_k(1+X_k+Y_k)$ choices of integers
$0\leq s_k\leq X_k+Y_k$.  Hence
\[
A(i,j)=\#(i\mathcal P_j\cap\setZ^{3j}).
\]
The defining matrix of $\mathcal P_j$ is totally unimodular: after changing
row signs, its nontrivial columns form a directed node-edge incidence matrix.
Thus $\mathcal P_j$ is a $3j$-dimensional lattice polytope, and its
$h^*$-polynomial is
\[
N_{2j}(t)\coloneqq\sum_{k=0}^{2j}\binom{2j}{k}^2t^k.
\]
In particular, its normalized volume is
$N_{2j}(1)=\binom{4j}{2j}$.

These numerator polynomials are transformed Legendre polynomials:
\[
N_m(t)=(1-t)^mP_m\!\left(\frac{1+t}{1-t}\right),
\]
where $P_m$ is the $m$th Legendre polynomial.  Their zeros are therefore
negative and simple.  Transporting the Legendre recurrence also gives
\[
(m+1)N_{m+1}(t)
=(2m+1)(1+t)N_m(t)-m(1-t)^2N_{m-1}(t).
\]
Thus every column $h^*$-vector is real-rooted, log-concave, and unimodal.

\begin{corollary}
For every $j\geq0$ there is a polynomial $p_j(i)$ of degree $j$ such that
\[
A(i,j)=\binom{i+j}{j}^2p_j(i).
\]
Its leading coefficient is
\[
\frac{\binom{4j}{2j}j!^2}{(3j)!},
\]
and it satisfies the reciprocity relation
\[
p_j(-i-j-1)=(-1)^jp_j(i).
\]
\end{corollary}

\begin{proof}
Regard the alternating formula for $A(i,j)$ as a polynomial in $i$.  For
$1\leq s\leq j$ and $0\leq m\leq j$, the product
\[
\binom{i-m+2j}{2j}
=\frac{\prod_{\ell=1}^{2j}(i-m+\ell)}{(2j)!}
\]
vanishes at $i=-s$, since $1\leq s+m\leq2j$.  Every summand therefore has a
double zero at each of $-1,-2,\dotsc,-j$.  This proves divisibility by
$\binom{i+j}{j}^2$.

The column generating function has numerator
$N_{2j}(1)=\binom{4j}{2j}\neq0$ at $t=1$.  Hence $A(i,j)$ has degree exactly
$3j$, with leading coefficient $\binom{4j}{2j}/(3j)!$.  Dividing by the
leading coefficient $1/j!^2$ of $\binom{i+j}{j}^2$ gives the stated degree
and leading coefficient of $p_j$.

Finally, $N_{2j}(t)$ is palindromic of degree $2j$.  Ehrhart reciprocity for
the $3j$-dimensional polytope $\mathcal P_j$, whose codegree is $j+1$, gives
\[
A(-i-j-1,j)=(-1)^{3j}A(i,j)=(-1)^jA(i,j).
\]
Since
$\binom{-i-1}{j}^2=\binom{i+j}{j}^2$, division gives the reciprocity relation
for $p_j$.
\end{proof}

The same calculation evaluates the bivariate generating function in the OEIS
entry.  Set
\[
\mathbf F(x,y)\coloneqq\sum_{i,j\geq0}A(i,j)x^iy^j,
\qquad
\xi=\frac{1+x}{1-x},\qquad w=\frac{y}{1-x}.
\]
Taking the even part of the Legendre generating function yields
\[
\mathbf F(x,y)=\frac{1}{2(1-x)}\left(
\frac{1}{\sqrt{1-2\xi\sqrt w+w}}
+\frac{1}{\sqrt{1+2\xi\sqrt w+w}}
\right).
\]
